%%%%%%%%%%%%%%%%%%%%%%%%%%%%%%%%%%%%%%%%%%%%%%%%%%%%%%%%%%%%%%%%%%%%%%
% Lớp 11 | Giới hạn. Hàm số liên tục | Hàm số liên tục
% Dạng 1D3?3-1: Câu hỏi lý thuyết
% Số câu: 9  (sinh tự động bởi tach_cau.py; từ nay thầy sửa trực tiếp file này)
%%%%%%%%%%%%%%%%%%%%%%%%%%%%%%%%%%%%%%%%%%%%%%%%%%%%%%%%%%%%%%%%%%%%%%

%%%%%%%%%%%%%%%% TRẮC NGHIỆM 4 PHƯƠNG ÁN %%%%%%%%%%%%%%%%

% nguồn: Ok_Toan11_TracNghiem_Lan2
\begin{ex}%[1D3N3-1]%[ma:77aaa64e]
	Cho hàm số $y=f(x)$ xác định trên  $K$, $x_0\in K$. Khi đó, hàm số $y=f(x)$ liên tục tại $x_0$ khi
	\choice 
		{$\lim\limits_{x \to  +\infty}f(x)=f\left( x_0 \right)$}
		{$\lim\limits_{x  \to -\infty}f(x)=f\left( x_0 \right)$}
		{$\lim\limits_{x \to x_0}f\left( x_0 \right)=f(x)$} 
		 {\True $\lim\limits_{x \to x_0}f(x)=f\left( x_0 \right)$}
	\loigiai{ 
		Hàm số $y=f(x)$ xác định trên $K$, $x_0\in K$. Hàm số $y=f(x)$ liên tục tại $x_0$ khi $\lim\limits_{x \to 3}f(x)=f\left(  x_0 \right)$.
	} 
\end{ex}

% nguồn: Ok_Toan11_TracNghiem_Lan2
\begin{ex}%[1D3N3-1]%[ma:7ad2523d]
	Cho hàm số $y=f(x)$ liên tục trên $\left(a;\,b \right)$. Điều kiện  cần và đủ để hàm số liên tục trên $\left[a;\,b \right]$ là
	\choice 
		 {$\lim\limits_{x \to a^+}f(x)=f\left( a \right)$ và $\lim\limits_{x \to b^+}f(x)=f\left( b \right)$}
		 {\True $\lim\limits_{x \to a^+}f(x)=f\left( a \right)$ và $\lim\limits_{x \to b^-}f(x)=f\left( b \right)$}
		{$\lim\limits_{x \to a^-}f(x)=f\left(  a \right)$ và $\lim\limits_{x \to b^+}f(x)=f\left( b \right)$}
		{$\lim\limits_{x \to a^-}f(x)=f\left( a \right)$  và $\lim\limits_{x \to b^-}f(x)=f\left( b \right)$}
	\loigiai{ 
		Hàm số $y=f(x)$ liên tục trên $\left(a;\,b \right)$. Điều kiện cần và đủ để hàm số liên tục trên $\left[a;\,b \right]$ là liên tục phải tại $a$ và liên  tục trái tại $b$, tức là $\lim\limits_{x \to a^+}f(x)=f\left( a \right)$ và $\lim\limits_{x \to b^-}f(x)=f\left( b \right)$.
	} 
\end{ex}

% nguồn: Ok_TracNghiemLop11_Lan1
\begin{ex}%[1D3N3-1]%[ma:7b5c4ef0]
	\immini{Hàm số $f(x)$ có đồ thị như hình bên. Hàm số gián đoạn tại điểm có hoành độ bằng  bao nhiêu?
	 \choice
		{$2$} 
		 {$3$}
		 {$0$}
		 {\True $1$}
		 }{
		\begin{tikzpicture}[>=stealth,line join=round,line cap=round,font=\footnotesize,scale=0.8]
				\draw[->]  (-2,0)--(4,0) node[below]{$x$};
				\draw[->]  (0,-0.5)--(0,4.5) node[left]{$y$};
				\fill[name=O] (0,0) circle (1pt) node[below left]  {$O$};
				\draw[smooth,samples=200,domain= -1.25:1]  plot (\x,{2*(\x)^2 +1});
				 \draw[smooth,samples=200,domain= 1:4] plot (\x,{\x-1});
				 \draw[dashed] (1,3)--(1,0)
					(1,3)--(0,3) 
					 (0,1)-|(2,0);
				 \coordinate[label=below:{$1$}] (1x) at (1,0);
				\coordinate[label=below left:{$1$}] (1y) at  (0,1);
				\coordinate[label=left:{$3$}]  (3y) at (0,3);
				\coordinate[label=below:{$2$}]  (2x) at (2,0);
				 \foreach \p in {1x,1y,2x,3y,{2,1}}
				\fill (\p)  circle (1.2pt);
			\end{tikzpicture}
		 }
	\loigiai{ 
		Ta có hàm số gián đoạn tại $x =1$. 
	} 
\end{ex}

% nguồn: Ok_TracNghiemLop11_Lan1
\begin{ex}%[1D3N3-1]%[ma:2ad7b494]
	Trong  các mệnh đề sau, mệnh đề nào \textbf{sai}?
	 \choice
		{Nếu hàm số $f(x)$ liên tục trên khoảng $(a;b)$ thì $f(x)$ cũng liên tục trên mọi khoảng con  của $(a;b)$}
		{Mọi hàm số đa thức đều liên  tục trên khoảng $(-\infty;+\infty)$}
		{Mọi hàm số phân thức hữu tỉ đều liên  tục trên mọi khoảng mà nó xác định}
		{\True Nếu hàm số $f(x)$ liên tục trên hai khoảng liên tiếp $(a;b)$  và $(b;c)$ thì $f(x)$ cũng liên tục trên khoảng $(a;c)$}
	\loigiai{ 
		Xét  hàm số
		$f(x)=\heva{&\sqrt{x-1}\text{  nếu }x\ge 1\\&2x+3\text{ nếu }x < 1.}$\\
	Ta thấy $f(x)$ liên tục trên các khoảng $(0;1)$ và $(1;2)$ nhưng không liên tục  trên $(0;2)$.}
\end{ex}

% nguồn: Ok_TracNghiemLop11_Lan1
\begin{ex}%[1D3N3-1]%[ma:1f9d28e2]
	Cho hàm số $y = f\left(x\right)$ xác định trên khoảng $K$ và $x_0 \in K$. Hàm số $y  = f\left(x\right)$ liên tục tại $x_0$ khi và chỉ khi
	\choice 
		{$\displaystyle\lim_{x \rightarrow x_0}  f\left(x\right) = +\infty$}
		{$\displaystyle\lim_{x \rightarrow x_0} f\left(x\right)  = -\infty$}
		{\True $\displaystyle\lim_{x  \rightarrow x_0} f\left(x\right) = f\left(x_0\right)$}
		 {$\displaystyle\lim_{x \rightarrow x_0^{+}} f\left(x\right) = f\left(x_0\right)$}
	\loigiai{ 
		Hàm số $y  = f\left(x\right)$ được gọi là liên tục tại $x_0$ nếu $\displaystyle\lim_{x \rightarrow x_0} f\left(x\right) = f\left(x_0 \right)$.
	} 
\end{ex}

% nguồn: Ok_TracNghiemLop11_Lan1
\begin{ex}%[1D3N3-1]%[ma:6f45fa67]
	Cho hàm số $y=f(x)$ xác định trên khoảng $(a;b)$ và $x_0 \in (a;b)$. Điều  kiện cần và đủ để hàm số $y=f(x)$ liên tục tại điểm $x_0$ là
	 \choice
		{$\displaystyle\lim\limits_{x \rightarrow x_0^+} f(x)  = f\left(x_0\right)$}
		{$\displaystyle\lim\limits_{x \rightarrow x_0^-} f(x) =  f\left(x_0\right)$}
		{$\displaystyle\lim\limits_{x \rightarrow x_0^+} f(x) = \displaystyle\lim\limits_{x \rightarrow x_0^-} f(x)$} 
		{\True $\displaystyle\lim\limits_{x \rightarrow x_0^+} f(x) = \displaystyle\lim\limits_{x \rightarrow x_0^-}  f(x) = f\left(x_0\right)$}
	 \loigiai{
		Điều kiện cần và đủ  để hàm số $y=f(x)$ liên tục tại điểm $x_0$ là
		\[\displaystyle\lim\limits_{x \rightarrow x_0^+} f(x) = \displaystyle\lim\limits_{x \rightarrow  x_0^-} f(x) = f\left(x_0\right).\]
	} 
\end{ex}

% nguồn: Ok_TracNghiemLop11_Lan1
\begin{ex}%[1D3N3-1]%[ma:b6e9bcd0]
	Trong các hàm số sau, hàm số nào  liên tục trên $\mathbb{R}$.
	\choice 
		{\True  $y=x^2+3x-4$}
		{$y=\dfrac{2x-2}{1-x}$} 
		{$y=\sqrt  x$}
		{$y=\tan  x$}
	\loigiai{ 
		\begin{itemize}
			\item Hàm số $y=x^2+3x-4$ là hàm số đa thức có tập xác định là $\mathbb{R}$ nên  liên tục trên $\mathbb{R}$.
			\item Các hàm số còn lại có tập  xác định không phải là $\mathbb{R}$ nên không liên tục trên $\mathbb{R}$.
		\end{itemize}
	 }
\end{ex}

% nguồn: Ok_TracNghiemLop11_Lan1
\begin{ex}%[1D3N3-1]%[ma:756929c9]
	Cho hàm số $f(x)=\dfrac{x+2}{x-4}$,  khi đó hàm số liên tục trên khoảng
	 \choice {$(-\infty; 5)$}{$(-3; 10)$}{$(-1;+\infty)$}{\True $(4;+\infty)$}
	\loigiai{Tập xác định của hàm số $f(x)$ là  $(-\infty; 4) \cup(4;+\infty)$.\\
		Vậy hàm số  liên tục trên các khoảng $(-\infty; 4)$ và $(4;+\infty)$.
	 }
\end{ex}

% nguồn: Ok_TracNghiemLop11_Lan1
\begin{ex}%[1D3N3-1]%[ma:cb53e36a]
	Cho hàm số $y = f(x)$ xác định trên $(a; b)$ và $x_0 \in (a;b)$. Trong các khẳng định sau, khẳng  định nào sau đây đúng?
	\choice 
		{Hàm số $y = f(x)$  liên tục tại $x_0$ nếu $f(x) = f(x_0)$}
		{\True Hàm số $y = f(x)$ được gọi là liên tục trên khoảng $[a; b]$ nếu  liên tục tại mọi điểm thuộc khoảng $(a; b)$ và $\lim\limits_{{x \to a^+}} f(x) = f(a), \lim\limits_{{x \to b^-}} f(x) = f(b)$}
		{Hàm số $y = f(x)$ được gọi là liên tục trên khoảng $[a; b]$ nếu liên tục tại mọi điểm thuộc khoảng $(a; b)$  và $\lim\limits_{{x \to a^-}} f(x) = f(a), \lim\limits_{{x \to b^+}} f(x) = f(b)$}
		{Hàm số $y = f(x)$ được gọi là liên  tục trên khoảng $(a; b)$ nếu liên tục tại điểm $x_0$ thuộc khoảng này}
	 \loigiai{
		Khẳng định \lq\lq Hàm số $y = f(x)$ được gọi là liên tục trên khoảng $[a; b]$ nếu liên tục tại mọi điểm thuộc khoảng $(a; b)$ và $\lim\limits_{{x \to a^+}} f(x)  = f(a), \lim\limits_{{x \to b^-}} f(x) = f(b)$\rq\rq \,này là khẳng định đúng.
	 }
\end{ex}
